\chapter{Sobolev Spaces}

\begin{remark}[Weak Poisson formulation]
\label{gt:rem-sobolev-weak-poisson}
\dcref{ch7:7.1--7.2}
\group{ch07-sobolev-spaces}
\source{gilbarg-trudinger:page-0156}
If $u\in C^2(\Omega)$ satisfies $\Delta u=f$, then
\begin{equation}
\tag{7.1}
\int_\Omega Du\cdot D\varphi\,dx=-\int_\Omega f\varphi\,dx
\qquad(\varphi\in C_0^1(\Omega)).
\end{equation}
The form
\begin{equation}
\tag{7.2}
(u,\varphi)=\int_\Omega Du\cdot D\varphi\,dx
\end{equation}
is an inner product on $C_0^1(\Omega)$. Its completion is
$W_0^{1,2}(\Omega)$. Whenever
$F(\varphi)=-\int_\Omega f\varphi\,dx$ extends boundedly to this completion,
the Riesz representation theorem produces $u\in W_0^{1,2}(\Omega)$ satisfying
$(u,\varphi)=F(\varphi)$.
\end{remark}

\section*{7.1. $L^p$ Spaces}
\label{gt:sec-lp-spaces}\dcref{ch7:7.1}\group{ch07-sobolev-spaces}\source{gilbarg-trudinger:page-0157}

Throughout the chapter, $\Omega\subset\mathbb R^n$ is a bounded domain.
Measurable functions are identified almost everywhere, and pointwise extrema
of such functions mean essential extrema.

\begin{definition}[$L^p$ spaces]
\label{gt:def-sobolev-lp-spaces}
\dcref{ch7:7.3--7.4}
\group{ch07-sobolev-spaces}
\source{gilbarg-trudinger:page-0157}
For $1\leq p<\infty$, let $L^p(\Omega)$ be the space of measurable functions
with finite norm
\begin{equation}
\tag{7.3}
\|u\|_{p;\Omega}=\|u\|_{L^p(\Omega)}
=\left(\int_\Omega |u|^p\,dx\right)^{1/p}.
\end{equation}
For vector- or matrix-valued $u$, $|u|$ is the Euclidean norm. Define
\begin{equation}
\tag{7.4}
\|u\|_{\infty;\Omega}=\|u\|_{L^\infty(\Omega)}
=\operatorname*{ess\,sup}_\Omega |u|.
\end{equation}
The domain subscript is omitted when it is clear.
\end{definition}

\begin{proposition}[Basic $L^p$ inequalities]
\label{gt:prop-sobolev-lp-inequalities}
\dcref{ch7:7.5--7.11}
\group{ch07-sobolev-spaces}
\source{gilbarg-trudinger:page-0157, gilbarg-trudinger:page-0158}
Let $a,b\geq0$ and let $p,q>1$ be conjugate exponents. Then
\begin{equation}
\tag{7.5}
ab\leq \frac{a^p}{p}+\frac{b^q}{q},
\end{equation}
and, for $\varepsilon>0$,
\begin{equation}
\tag{7.6}
ab\leq \frac{\varepsilon a^p}{p}
+\frac{\varepsilon^{-q/p}b^q}{q}
\leq \varepsilon a^p+\varepsilon^{-q/p}b^q.
\end{equation}
If $u\in L^p(\Omega)$ and $v\in L^q(\Omega)$, then
\begin{equation}
\tag{7.7}
\int_\Omega |uv|\,dx\leq \|u\|_p\|v\|_q.
\end{equation}
If $1\leq p\leq q\leq\infty$ and $u\in L^q(\Omega)$, then
\begin{equation}
\tag{7.8}
|\Omega|^{-1/p}\|u\|_p
\leq |\Omega|^{-1/q}\|u\|_q,
\end{equation}
If moreover $q\leq r\leq\infty$, $u\in L^r(\Omega)$, and
$q^{-1}=\lambda p^{-1}+(1-\lambda)r^{-1}$, then
\begin{equation}
\tag{7.9}
\|u\|_q\leq \|u\|_p^\lambda\|u\|_r^{1-\lambda}.
\end{equation}
Consequently, for $p<q<r$ and $\varepsilon>0$,
\begin{equation}
\tag{7.10}
\|u\|_q\leq \varepsilon\|u\|_r
+\varepsilon^{-\mu}\|u\|_p,
\qquad
\mu=\frac{p^{-1}-q^{-1}}{q^{-1}-r^{-1}}.
\end{equation}
More generally, if $u_j\in L^{p_j}(\Omega)$ and
$\sum_{j=1}^m p_j^{-1}=1$, then
\begin{equation}
\tag{7.11}
\int_\Omega |u_1\cdots u_m|\,dx
\leq \prod_{j=1}^m\|u_j\|_{p_j}.
\end{equation}
\end{proposition}
\begin{proof}
\sketch Young's inequality gives (7.5), and rescaling its two arguments gives
(7.6). Integrating the pointwise inequality after normalization proves
H\"older's inequality. Apply it to $u$ and $1$ for (7.8), and to suitable
powers of $u$ for (7.9); (7.10) then follows from (7.6). Induction on $m$
proves (7.11).
\end{proof}

\begin{definition}[Normalized $L^p$ functional]
\label{gt:def-sobolev-normalized-lp-functional}
\dcref{ch7:7.12}
\group{ch07-sobolev-spaces}
\source{gilbarg-trudinger:page-0158}
For $p>0$, set
\begin{equation}
\tag{7.12}
\Phi_p(u)=\left(\frac1{|\Omega|}\int_\Omega |u|^p\,dx\right)^{1/p}.
\end{equation}
\end{definition}

\begin{proposition}[Dependence of $\Phi_p$ on $p$]
\label{gt:prop-sobolev-normalized-lp-functional}
\dcref{ch7:7.8--7.10}
\group{ch07-sobolev-spaces}
\source{gilbarg-trudinger:page-0158}
For fixed $u$, $\Phi_p(u)$ is nondecreasing in $p$ and logarithmically convex
as a function of $p^{-1}$. For $p\geq1$,
$\Phi_p(u)=|\Omega|^{-1/p}\|u\|_p$.
\end{proposition}
\begin{proof}
\sketch The monotonicity is (7.8) after normalization by $|\Omega|$;
logarithmic convexity is (7.9).
\end{proof}

For $p<\infty$, $L^p(\Omega)$ is separable. Its dual is $L^{p'}(\Omega)$,
where $p^{-1}+(p')^{-1}=1$, and it is reflexive when $1<p<\infty$.
The space $L^2(\Omega)$ has inner product
$ (u,v)=\int_\Omega uv\,dx$.

\section*{7.2. Regularization and Smooth Approximation}

Convergence in $L^p_{\mathrm{loc}}(\Omega)$ means convergence in
$L^p(\Omega')$ for every $\Omega'\Subset\Omega$. Fix a nonnegative
$\rho\in C_0^\infty(B_1(0))$ with $\int\rho=1$.

\begin{definition}[Mollification]
\label{gt:def-sobolev-mollification}
\dcref{ch7:7.13}
\group{ch07-sobolev-spaces}
\source{gilbarg-trudinger:page-0159}
For $u\in L^1_{\mathrm{loc}}(\Omega)$ and
$0<h<\dist(x,\partial\Omega)$, define
\begin{equation}
\tag{7.13}
u_h(x)=h^{-n}\int_\Omega
\rho\!\left(\frac{x-y}{h}\right)u(y)\,dy.
\end{equation}
Then $u_h$ is smooth wherever it is defined. If $u\in L^1(\Omega)$ is
extended by zero, the same formula defines a smooth function on $\mathbb R^n$.
\end{definition}

\begin{lemma}[Uniform convergence of mollifications]
\label{gt:lem-sobolev-regularization-uniform}\dcref{ch7:7.1}\group{ch07-sobolev-spaces}
\source{gilbarg-trudinger:page-0159}
If $u\in C^0(\Omega)$, then $u_h\to u$ uniformly on every
$\Omega'\Subset\Omega$.
\end{lemma}
\begin{proof}
\sketch The identity
$u_h(x)=\int_{B_1}\rho(z)u(x-hz)\,dz$ bounds
$|u_h(x)-u(x)|$ by the modulus of continuity of $u$ on a fixed compact
neighborhood of $\Omega'$, which tends to zero with $h$.
\end{proof}

If $u\in C^\alpha(\Omega)$ and
$\Omega''=\{x:\dist(x,\Omega')<h\}\Subset\Omega$, then
\begin{equation}
\tag{7.14}
[u_h]_{\alpha;\Omega'}\leq [u]_{\alpha;\Omega''}.
\end{equation}
Thus $u_h\to u$ in $C^{\alpha'}(\Omega')$ whenever
$0\leq\alpha'<\alpha$.

\begin{lemma}[$L^p$ convergence of mollifications]
\label{gt:lem-sobolev-regularization-lp}\dcref{ch7:7.2}\group{ch07-sobolev-spaces}
\source{gilbarg-trudinger:page-0160, gilbarg-trudinger:page-0161}
Let $1\leq p<\infty$. If $u\in L^p_{\mathrm{loc}}(\Omega)$, then
$u_h\to u$ in $L^p_{\mathrm{loc}}(\Omega)$. If $u\in L^p(\Omega)$ is
extended by zero, convergence holds in $L^p(\mathbb R^n)$.
Moreover, for $\Omega'\Subset\Omega$ and
$\Omega''=\{x:\dist(x,\Omega')<h\}\Subset\Omega$,
\begin{equation}
\tag{7.15}
\|u_h\|_{L^p(\Omega')}\leq\|u\|_{L^p(\Omega'')}.
\end{equation}
\end{lemma}
\begin{proof}
\sketch Jensen's inequality and Fubini give (7.15). Approximate $u$ on a
slightly larger subdomain by a continuous function, apply Lemma~7.1 to that
function, and use (7.15) on the approximation error. The global claim follows
after zero extension.
\end{proof}

\section*{7.3. Weak Derivatives}

\begin{definition}[Weak derivative]
\label{gt:def-sobolev-weak-derivative}
\dcref{ch7:7.16}
\group{ch07-sobolev-spaces}
\source{gilbarg-trudinger:page-0161, gilbarg-trudinger:page-0162}
Let $u\in L^1_{\mathrm{loc}}(\Omega)$ and let $\alpha$ be a multi-index.
A function $v\in L^1_{\mathrm{loc}}(\Omega)$ is the weak derivative
$D^\alpha u$ if
\begin{equation}
\tag{7.16}
\int_\Omega \varphi v\,dx
=(-1)^{|\alpha|}\int_\Omega uD^\alpha\varphi\,dx
\qquad(\varphi\in C_0^{|\alpha|}(\Omega)).
\end{equation}
It is unique almost everywhere. Write $W^k(\Omega)$ for the functions whose
weak derivatives through order $k$ exist locally.
\end{definition}

\begin{lemma}[Weak derivatives commute with mollification]
\label{gt:lem-sobolev-mollifier-derivative}\dcref{ch7:7.3}\group{ch07-sobolev-spaces}
\source{gilbarg-trudinger:page-0162}
If $u\in L^1_{\mathrm{loc}}(\Omega)$ and $D^\alpha u$ exists, then, whenever
$\dist(x,\partial\Omega)>h$,
\begin{equation}
\tag{7.17}
D^\alpha u_h(x)=(D^\alpha u)_h(x).
\end{equation}
\end{lemma}
\begin{proof}
\sketch Differentiate the convolution kernel in $x$, replace its
$x$-derivative by the signed $y$-derivative, and apply (7.16).
\end{proof}

\begin{theorem}[Weak and strong derivatives]
\label{gt:thm-sobolev-weak-derivative-approx}\dcref{ch7:7.4}\group{ch07-sobolev-spaces}
\source{gilbarg-trudinger:page-0162, gilbarg-trudinger:page-0163}
For $u,v\in L^1_{\mathrm{loc}}(\Omega)$, one has $v=D^\alpha u$ if and only
if there are $u_m\in C^\infty(\Omega)$ such that
$u_m\to u$ and $D^\alpha u_m\to v$ in $L^1_{\mathrm{loc}}(\Omega)$.
\end{theorem}
\begin{proof}
\sketch If $v=D^\alpha u$, mollify on a compact exhaustion and choose a
diagonal sequence, using Lemmas~7.2 and 7.3. Conversely, pass to the limit in
the classical integration-by-parts identity for $u_m$.
\end{proof}

\begin{proposition}[Calculus for weak derivatives]
\label{gt:prop-sobolev-weak-calculus}
\dcref{ch7:7.18--7.19}
\group{ch07-sobolev-spaces}
\source{gilbarg-trudinger:page-0162, gilbarg-trudinger:page-0163}
If $u,v\in W^1(\Omega)$ and both sides below are locally integrable, then
\begin{equation}
\tag{7.18}
D(uv)=uDv+vDu.
\end{equation}
If $\psi:\Omega\to\widetilde\Omega$ is a $C^1$ diffeomorphism,
$u\in W^1(\Omega)$, $v=u\circ\psi^{-1}$, and $y=\psi(x)$, then
\begin{equation}
\tag{7.19}
D_i u(x)=\frac{\partial y_j}{\partial x_i}D_{y_j}v(y)
\quad\text{for almost every }x\in\Omega.
\end{equation}
Every locally Lipschitz function is weakly differentiable. More generally,
$u$ is weakly differentiable exactly when it has a representative absolutely
continuous on almost every coordinate-parallel line and its classical partial
derivatives are locally integrable.
\end{proposition}

\section*{7.4. The Chain Rule}

\begin{lemma}[Chain rule for a $C^1$ function]
\label{gt:lem-sobolev-chain-rule-c1}\dcref{ch7:7.5}\group{ch07-sobolev-spaces}
\source{gilbarg-trudinger:page-0163, gilbarg-trudinger:page-0164}
If $f\in C^1(\mathbb R)$, $f'\in L^\infty(\mathbb R)$, and
$u\in W^1(\Omega)$, then $f\circ u\in W^1(\Omega)$ and
$D(f\circ u)=f'(u)Du$.
\end{lemma}
\begin{proof}
\sketch Choose smooth $u_m$ converging to $u$ together with their gradients
locally in $L^1$. Boundedness of $f'$ gives $f(u_m)\to f(u)$ in $L^1$;
after passage to an almost-everywhere convergent subsequence, continuity of
$f'$ and dominated convergence give $f'(u_m)Du_m\to f'(u)Du$ locally in
$L^1$. Apply Theorem~7.4.
\end{proof}

For a real-valued $u$, set $u^+=\max\{u,0\}$ and
$u^-=\min\{u,0\}$.

\begin{lemma}[Positive and negative parts]
\label{gt:lem-sobolev-positive-negative-parts}\dcref{ch7:7.6}\group{ch07-sobolev-spaces}
\source{gilbarg-trudinger:page-0164}
If $u\in W^1(\Omega)$, then $u^+,u^-,|u|\in W^1(\Omega)$ and
\begin{equation}
\tag{7.20}
Du^+=\begin{cases}Du&u>0,\\0&u\leq0,
\end{cases}
\quad
Du^-=\begin{cases}0&u\geq0,\\Du&u<0,
\end{cases}
\quad
D|u|=\begin{cases}Du&u>0,\\0&u=0,\\-Du&u<0.
\end{cases}
\end{equation}
\end{lemma}
\begin{proof}
\sketch Apply Lemma~7.5 to smooth bounded-slope approximations of the positive
part and pass to the limit in (7.16). Obtain the remaining formulas from
$u^-=-(-u)^+$ and $|u|=u^+-u^-$.
\end{proof}

\begin{lemma}[Gradient on a level set]
\label{gt:lem-sobolev-constant-gradient-zero}\dcref{ch7:7.7}\group{ch07-sobolev-spaces}
\source{gilbarg-trudinger:page-0164, gilbarg-trudinger:page-0165}
If $u\in W^1(\Omega)$, then $Du=0$ almost everywhere on every measurable set
on which $u$ is constant.
\end{lemma}
\begin{proof}
\sketch Subtract the constant and use $Du=Du^++Du^-$ together with (7.20).
\end{proof}

\begin{theorem}[Piecewise smooth chain rule]
\label{gt:thm-sobolev-piecewise-chain-rule}\dcref{ch7:7.8}\group{ch07-sobolev-spaces}
\source{gilbarg-trudinger:page-0165}
Let $f:\mathbb R\to\mathbb R$ be continuous and piecewise $C^1$, with
$f'\in L^\infty(\mathbb R)$, and let $L$ be its set of corner points. If
$u\in W^1(\Omega)$, then $f\circ u\in W^1(\Omega)$ and
\begin{equation}
\tag{7.21}
D(f\circ u)=
\begin{cases}
f'(u)Du,&u\notin L,\\
0,&u\in L.
\end{cases}
\end{equation}
\end{theorem}
\begin{proof}
\sketch Reduce to one corner, translate it to zero, and represent $f(u)$ using
two $C^1$ bounded-slope extensions evaluated at $u^+$ and $u^-$. Apply
Lemmas~7.5--7.7 and iterate over the corners.
\end{proof}

The same formula holds for a Lipschitz $f$ whenever a measurable choice
$h=f'$ off the corner set satisfies $h(u)Du\in L^1_{\mathrm{loc}}(\Omega)$;
on the corner level sets, Lemma~7.7 makes the choice of $h$ immaterial.

\section*{7.5. The $W^{k,p}$ Spaces}
\label{gt:sec-wkp-spaces}\dcref{ch7:7.5}\group{ch07-sobolev-spaces}\source{gilbarg-trudinger:page-0165}

\begin{definition}[Sobolev spaces]
\label{gt:def-sobolev-wkp-spaces}
\dcref{ch7:7.22--7.23}
\group{ch07-sobolev-spaces}
\source{gilbarg-trudinger:page-0165, gilbarg-trudinger:page-0166}
For an integer $k\geq0$ and $1\leq p<\infty$, define
\[
W^{k,p}(\Omega)=\{u\in W^k(\Omega):D^\alpha u\in L^p(\Omega)
\text{ for }|\alpha|\leq k\}
\]
with norm
\begin{equation}
\tag{7.22}
\|u\|_{k,p;\Omega}
=\left(\int_\Omega\sum_{|\alpha|\leq k}|D^\alpha u|^p\,dx\right)^{1/p}.
\end{equation}
An equivalent norm is
\begin{equation}
\tag{7.23}
\|u\|_{W^{k,p}(\Omega)}=\sum_{|\alpha|\leq k}\|D^\alpha u\|_p.
\end{equation}
The space $W_0^{k,p}(\Omega)$ is the closure of $C_0^k(\Omega)$ in
$W^{k,p}(\Omega)$, and $W^{k,p}_{\mathrm{loc}}(\Omega)$ consists of functions
belonging to $W^{k,p}(\Omega')$ for every $\Omega'\Subset\Omega$.
\end{definition}

Both $W^{k,p}(\Omega)$ and $W_0^{k,p}(\Omega)$ are Banach spaces, separable
for $p<\infty$ and reflexive for $1<p<\infty$. For $p=2$ they are Hilbert
spaces under
\begin{equation}
\tag{7.24}
(u,v)_k=\int_\Omega\sum_{|\alpha|\leq k}D^\alpha uD^\alpha v\,dx.
\end{equation}
The chain rule (7.21) holds in $W^{1,p}$, and it preserves $W_0^{1,p}$ when
$f(0)=0$. A compactly supported member of $W^{k,p}_{\mathrm{loc}}(\Omega)$
lies in $W_0^{k,p}(\Omega)$. Moreover,
$W^{k,\infty}_{\mathrm{loc}}(\Omega)=C^{k-1,1}(\Omega)$; for Lipschitz
$\Omega$, $W^{k,\infty}(\Omega)=C^{k-1,1}(\overline\Omega)$.

\section*{7.6. Density Theorems}

\begin{theorem}[Interior smooth density]
\label{gt:thm-sobolev-density-smooth}\dcref{ch7:7.9}\group{ch07-sobolev-spaces}
\source{gilbarg-trudinger:page-0166, gilbarg-trudinger:page-0167}
The subspace $C^\infty(\Omega)\cap W^{k,p}(\Omega)$ is dense in
$W^{k,p}(\Omega)$.
\end{theorem}
\begin{proof}
\sketch Choose an exhaustion $\Omega_j\Subset\Omega_{j+1}$ and a subordinate
partition of unity $\{\psi_j\}$. For any $u\in W^{k,p}(\Omega)$ and
$\varepsilon>0$, choose $h_j$ so that
\begin{equation}
\tag{7.25}
h_j\leq\dist(\Omega_j,\partial\Omega_{j+1})
\end{equation}
and
$\|(\psi_ju)_{h_j}-\psi_ju\|_{W^{k,p}(\Omega)}\leq\varepsilon 2^{-j}$.
The locally finite sum $v=\sum_j(\psi_ju)_{h_j}$ is smooth and satisfies
$\|u-v\|_{W^{k,p}(\Omega)}\leq\varepsilon$.
\end{proof}

\begin{proposition}[Boundary smooth density under a segment condition]
\label{gt:prop-sobolev-segment-density}
\dcref{ch7:7.9}
\group{ch07-sobolev-spaces}
\source{gilbarg-trudinger:page-0167}
Suppose $\partial\Omega$ has a locally finite open cover $\{U_i\}$ and vectors
$y^i$ such that $x+ty^i\in\Omega$ for every
$x\in\overline\Omega\cap U_i$ and $0<t<1$. Then
$C^\infty(\overline\Omega)$ is dense in $W^{k,p}(\Omega)$.
\end{proposition}

\section*{7.7. Embedding Theorems}

\begin{theorem}[First-order Sobolev embeddings]
\label{gt:thm-sobolev-embedding-basic}\dcref{ch7:7.10}\group{ch07-sobolev-spaces}
\source{gilbarg-trudinger:page-0167, gilbarg-trudinger:page-0168, gilbarg-trudinger:page-0169, gilbarg-trudinger:page-0170}
\source{gilbarg-trudinger:page-0168}
\source{gilbarg-trudinger:page-0169}
\source{gilbarg-trudinger:page-0170}
Let $u\in W_0^{1,p}(\Omega)$. If $1\leq p<n$, then
$u\in L^{np/(n-p)}(\Omega)$ and
\begin{equation}
\tag{7.26}
\|u\|_{np/(n-p)}\leq C\|Du\|_p.
\end{equation}
If $p>n$, then $u\in C^0(\overline\Omega)$ and
\[
\sup_\Omega|u|\leq C|\Omega|^{1/n-1/p}\|Du\|_p.
\]
In both cases $C=C(n,p)$.
\end{theorem}
\begin{proof}
\sketch For $u\in C_0^1(\Omega)$, integration along each coordinate line
gives
\begin{equation}
\tag{7.27}
|u(x)|^{n/(n-1)}
\leq\left(\prod_{i=1}^n\int_{-\infty}^{\infty}|D_i u|\,dx_i
\right)^{1/(n-1)}.
\end{equation}
Successive integration and generalized H\"older yield
\begin{equation}
\tag{7.28}
\|u\|_{n/(n-1)}\leq n^{-1/2}\|Du\|_1.
\end{equation}
Apply (7.28) to $|u|^\gamma$, use H\"older's inequality, and choose
$\gamma=(n-1)p/(n-p)$ to obtain (7.26). For $p>n$, iteration of the same
estimate and scaling give the stated supremum bound. Approximation by
$C_0^1(\Omega)$ extends both estimates to $W_0^{1,p}(\Omega)$.
\end{proof}

\begin{remark}[Sharp constant in (7.26)]
\label{gt:rem-sobolev-sharp-constant}
\dcref{ch7:7.26}
\group{ch07-sobolev-spaces}
\source{gilbarg-trudinger:page-0170}
For $1<p<n$, the optimal constant in (7.26) is
\[
C=\frac1{n\sqrt\pi}
\left(\frac{n!\,\Gamma(n/2)}{2\Gamma(n/p)\Gamma(n+1-n/p)}\right)^{1/n}
\gamma^{1-1/p},
\qquad \gamma=\frac{n(p-1)}{n-p}.
\]
At $p=1$ the limiting value is $n^{-1}\omega_n^{-1/n}$.
\end{remark}

\begin{corollary}[Higher-order Sobolev embeddings]
\label{gt:cor-sobolev-embedding-higher-order}\dcref{ch7:7.11}\group{ch07-sobolev-spaces}
\source{gilbarg-trudinger:page-0170}
There are continuous embeddings
\[
W_0^{k,p}(\Omega)\hookrightarrow
\begin{cases}
L^{np/(n-kp)}(\Omega),&kp<n,\\
C^m(\overline\Omega),&0\leq m<k-n/p.
\end{cases}
\]
On $W_0^{k,p}(\Omega)$, an equivalent norm is
\begin{equation}
\tag{7.29}
\|u\|_{W_0^{k,p}(\Omega)}
=\left(\int_\Omega\sum_{|\alpha|=k}|D^\alpha u|^p\,dx\right)^{1/p}.
\end{equation}
\end{corollary}
\begin{proof}
\sketch Iterate Theorem~7.10 through derivatives of orders below $k$. The
same inequalities control all lower-order terms by the top-order expression
in (7.29).
\end{proof}

\begin{proposition}[Embeddings under a uniform cone condition]
\label{gt:prop-sobolev-cone-embedding}
\dcref{ch7:7.10}
\group{ch07-sobolev-spaces}
\source{gilbarg-trudinger:page-0170}
Assume there is a fixed cone $K_\Omega$ such that every $x\in\Omega$ is the
vertex of a congruent cone $K_\Omega(x)\subset\overline\Omega$. Then
\[
W^{k,p}(\Omega)\hookrightarrow
\begin{cases}
L^{np/(n-kp)}(\Omega),&kp<n,\\
C_B^m(\Omega),&0\leq m<k-n/p,
\end{cases}
\]
where
$C_B^m(\Omega)=\{u\in C^m(\Omega):D^\alpha u\in L^\infty(\Omega)
\text{ for }|\alpha|\leq m\}$.
\end{proposition}

\section*{7.8. Potential Estimates and Embedding Theorems}
\label{gt:sec-potential-estimates}\dcref{ch7:7.8}\group{ch07-sobolev-spaces}\source{gilbarg-trudinger:page-0171}

\begin{definition}[Riesz potential]
\label{gt:def-sobolev-riesz-potential}
\dcref{ch7:7.31}
\group{ch07-sobolev-spaces}
\source{gilbarg-trudinger:page-0171}
For $0<\mu\leq1$ and $f\in L^1(\Omega)$, define
\begin{equation}
\tag{7.31}
(V_\mu f)(x)=\int_\Omega |x-y|^{n(\mu-1)}f(y)\,dy.
\end{equation}
\end{definition}

\begin{proposition}[Potential of the constant function]
\label{gt:prop-sobolev-riesz-potential-one}
\dcref{ch7:7.32}
\group{ch07-sobolev-spaces}
\source{gilbarg-trudinger:page-0171}
For $0<\mu\leq1$,
\begin{equation}
\tag{7.32}
V_\mu1\leq \mu^{-1}\omega_n^{1-\mu}|\Omega|^\mu.
\end{equation}
\end{proposition}
\begin{proof}
\sketch Choose $R$ with $|B_R|=|\Omega|$. Rearrangement by distance from
$x$ bounds the integral over $\Omega$ by the integral over $B_R(x)$, whose
polar-coordinate value is the right side of (7.32).
\end{proof}

\begin{lemma}[Mapping estimate for Riesz potentials]
\label{gt:lem-sobolev-riemann-potential-mapping}\dcref{ch7:7.12}\group{ch07-sobolev-spaces}
\source{gilbarg-trudinger:page-0171, gilbarg-trudinger:page-0172}
\source{gilbarg-trudinger:page-0172}
Let $1\leq p,q\leq\infty$ satisfy
\begin{equation}
\tag{7.33}
0\leq\delta=p^{-1}-q^{-1}<\mu.
\end{equation}
Then $V_\mu:L^p(\Omega)\to L^q(\Omega)$ is bounded and
\begin{equation}
\tag{7.34}
\|V_\mu f\|_q
\leq\left(\frac{1-\delta}{\mu-\delta}\right)^{1-\delta}
\omega_n^{1-\mu}|\Omega|^{\mu-\delta}\|f\|_p.
\end{equation}
\end{lemma}
\begin{proof}
\sketch Put $r^{-1}=1+q^{-1}-p^{-1}=1-\delta$. Estimate the
$L^r$ norm of $|x-y|^{n(\mu-1)}$ using (7.32), then apply the generalized
H\"older inequality in the standard convolution estimate.
\end{proof}

\begin{remark}[Endpoint potential mapping]
\label{gt:rem-sobolev-riesz-endpoint}
\dcref{ch7:7.12}
\group{ch07-sobolev-spaces}
\source{gilbarg-trudinger:page-0172}
For $p>1$, the mapping conclusion of Lemma~7.12 remains valid at
$\delta=\mu$. In particular, $V_\mu:L^p(\Omega)\to L^\infty(\Omega)$ is
bounded when $p>\mu^{-1}$.
\end{remark}

\begin{lemma}[Critical exponential estimate]
\label{gt:lem-sobolev-riemann-potential-exponential}\dcref{ch7:7.13}\group{ch07-sobolev-spaces}
\source{gilbarg-trudinger:page-0172, gilbarg-trudinger:page-0173}
\source{gilbarg-trudinger:page-0173}
Let $1<p<\infty$, $f\in L^p(\Omega)$, and $g=V_{1/p}f$. There are
$c_1,c_2>0$, depending only on $n,p$, such that
\begin{equation}
\tag{7.35}
\int_\Omega
\exp\!\left[\left(\frac{|g|}{c_1\|f\|_p}\right)^{p'}\right]dx
\leq c_2|\Omega|,
\qquad p'=\frac p{p-1}.
\end{equation}
\end{lemma}
\begin{proof}
\sketch Apply (7.34) with arbitrary $q\geq p$ to bound the $q$th moments
of $g$. Insert the resulting bounds into the exponential series. Stirling's
estimate makes the series converge once $c_1^{p'}>e\omega_np'$.
\end{proof}

\begin{lemma}[Newtonian representation]
\label{gt:lem-sobolev-newtonian-representation}\dcref{ch7:7.14}\group{ch07-sobolev-spaces}
\source{gilbarg-trudinger:page-0173, gilbarg-trudinger:page-0174}
\source{gilbarg-trudinger:page-0174}
If $u\in W_0^{1,1}(\Omega)$, then
\begin{equation}
\tag{7.36}
u(x)=\frac1{n\omega_n}\int_\Omega
\frac{(x_i-y_i)D_i u(y)}{|x-y|^n}\,dy
\quad\text{for almost every }x\in\Omega,
\end{equation}
and hence
\begin{equation}
\tag{7.37}
|u|\leq\frac1{n\omega_n}V_{1/n}|Du|.
\end{equation}
\end{lemma}
\begin{proof}
\sketch For $u\in C_0^1(\Omega)$, integrate the directional derivative of
$u$ along every ray from $x$ and average over the unit sphere. This gives
(7.36); density and Lemma~7.12 pass the identity to $W_0^{1,1}(\Omega)$.
Taking absolute values yields (7.37).
\end{proof}

\begin{theorem}[Moser--Trudinger estimate]
\label{gt:thm-sobolev-moser-trudinger}\dcref{ch7:7.15}\group{ch07-sobolev-spaces}
\source{gilbarg-trudinger:page-0174}
If $u\in W_0^{1,n}(\Omega)$, there are $c_1,c_2>0$, depending only on $n$,
such that
\begin{equation}
\tag{7.38}
\int_\Omega
\exp\!\left[\left(\frac{|u|}{c_1\|Du\|_n}\right)^{n/(n-1)}\right]dx
\leq c_2|\Omega|.
\end{equation}
\end{theorem}
\begin{proof}
\sketch Combine the pointwise estimate (7.37) with Lemma~7.13 for $p=n$.
\end{proof}

\begin{corollary}[Higher-order critical exponential estimate]
\label{gt:cor-sobolev-moser-trudinger-higher-order}
\dcref{ch7:7.15}
\group{ch07-sobolev-spaces}
\source{gilbarg-trudinger:page-0174}
For $u\in W_0^{k,1}(\Omega)$,
\begin{equation}
\tag{7.39}
|u|\leq\frac1{(k-1)!\,n\omega_n}V_{k/n}|D^ku|.
\end{equation}
If $n=kp$ and $u\in W_0^{k,p}(\Omega)$, there are $c_1,c_2>0$, depending
only on $n,k$, such that
\begin{equation}
\tag{7.40}
\int_\Omega
\exp\!\left[\left(\frac{|u|}{c_1\|D^ku\|_p}\right)^{p/(p-1)}\right]dx
\leq c_2|\Omega|.
\end{equation}
\end{corollary}
\begin{proof}
\sketch Iterate the ray representation used for (7.36) to obtain (7.39),
then apply Lemma~7.13 with $\mu=k/n=1/p$.
\end{proof}

\begin{lemma}[Convex-domain average estimate]
\label{gt:lem-sobolev-convex-average-estimate}\dcref{ch7:7.16}\group{ch07-sobolev-spaces}
\source{gilbarg-trudinger:page-0174, gilbarg-trudinger:page-0175}
\source{gilbarg-trudinger:page-0175}
Let $\Omega$ be convex, $u\in W^{1,1}(\Omega)$, and let
$S\subset\Omega$ be measurable with $|S|>0$. If
$u_S=|S|^{-1}\int_Su\,dx$ and $d=\diam\Omega$, then
\begin{equation}
\tag{7.41}
|u(x)-u_S|\leq\frac{d^n}{n|S|}
\int_\Omega |x-y|^{1-n}|Du(y)|\,dy
\quad\text{for almost every }x\in\Omega.
\end{equation}
\end{lemma}
\begin{proof}
\sketch First take $u\in C^1(\Omega)$. Integrate $Du$ along the segment from
$x$ to each $y\in S$, average over $S$, and pass to polar coordinates about
$x$; convexity keeps every segment in $\Omega$ and gives the factor
$d^n/(n|S|)$. Apply Theorem~7.9 to general $u$.
\end{proof}

\begin{theorem}[Morrey embedding]
\label{gt:thm-sobolev-morrey}\dcref{ch7:7.17}\group{ch07-sobolev-spaces}
\source{gilbarg-trudinger:page-0175, gilbarg-trudinger:page-0176}
\source{gilbarg-trudinger:page-0176}
Let $p>n$, $\gamma=1-n/p$, and $u\in W_0^{1,p}(\Omega)$. Then
$u\in C^\gamma(\overline\Omega)$ and, for every ball $B_R$,
\begin{equation}
\tag{7.42}
\operatorname*{osc}_{\Omega\cap B_R}u
\leq C R^\gamma\|Du\|_p,
\qquad C=C(n,p).
\end{equation}
\end{theorem}
\begin{proof}
\sketch Apply (7.41) and (7.34) on a ball, with
$S=\Omega=B_R$, $q=\infty$, and $\mu=1/n$. The resulting estimate for
$|u-u_B|$ at two points gives (7.42).
\end{proof}

\begin{corollary}[Global Sobolev and Poincare estimates]
\label{gt:cor-sobolev-global-estimates}
\dcref{ch7:7.17}
\group{ch07-sobolev-spaces}
\source{gilbarg-trudinger:page-0176}
For $p>n$, $\gamma=1-n/p$, and $u\in W_0^{1,p}(\Omega)$,
\begin{equation}
\tag{7.43}
|u|_{0,\gamma}\leq C\bigl[1+(\diam\Omega)^\gamma\bigr]\|Du\|_p.
\end{equation}
At the critical exponent $p=n$, $W_0^{1,n}(\Omega)$ embeds into the Orlicz
space associated with
$\varphi(t)=\exp(|t|^{n/(n-1)})-1$.
For $1\leq p<\infty$ and $u\in W_0^{1,p}(\Omega)$,
\begin{equation}
\tag{7.44}
\|u\|_p\leq\left(\frac{|\Omega|}{\omega_n}\right)^{1/n}\|Du\|_p.
\end{equation}
If $\Omega$ is convex, $u\in W^{1,p}(\Omega)$, and
$S\subset\Omega$ is measurable with $|S|>0$, then
\begin{equation}
\tag{7.45}
\|u-u_S\|_p\leq
\left(\frac{\omega_n}{|S|}\right)^{1-1/n}
(\diam\Omega)^n\|Du\|_p.
\end{equation}
\end{corollary}
\begin{proof}
\sketch Equation (7.43) combines Theorems~7.10 and 7.17. The critical
embedding is (7.38). Equations (7.44) and (7.45) follow respectively by
combining Lemma~7.12 with (7.37) and (7.41).
\end{proof}

\section*{7.9. Morrey and John--Nirenberg Estimates}
\label{gt:sec-morrey-john-nirenberg}\dcref{ch7:7.9}\group{ch07-sobolev-spaces}\source{gilbarg-trudinger:page-0176}

\begin{definition}[Morrey spaces]
\label{gt:def-sobolev-morrey-space}
\dcref{ch7:7.46}
\group{ch07-sobolev-spaces}
\source{gilbarg-trudinger:page-0176}
For $1\leq p\leq\infty$, an integrable function $f$ belongs to
$M^p(\Omega)$ if some $K$ satisfies
\begin{equation}
\tag{7.46}
\int_{\Omega\cap B_R}|f|\,dx\leq KR^{n(1-1/p)}
\end{equation}
for every ball $B_R$. The norm $\|f\|_{M^p(\Omega)}$ is the infimum of such
$K$.
\end{definition}

\begin{proposition}[Elementary Morrey inclusions]
\label{gt:prop-sobolev-morrey-inclusions}
\dcref{ch7:7.46}
\group{ch07-sobolev-spaces}
\source{gilbarg-trudinger:page-0176}
For $1\leq p\leq\infty$, one has $L^p(\Omega)\subset M^p(\Omega)$.
Moreover, $M^1(\Omega)=L^1(\Omega)$ and
$M^\infty(\Omega)=L^\infty(\Omega)$.
\end{proposition}
\begin{proof}
\sketch Apply H\"older's inequality on $\Omega\cap B_R$. At $p=1$ the
defining bound is the $L^1$ bound, while at $p=\infty$ differentiation of
integrals over shrinking balls identifies the essential supremum.
\end{proof}

\begin{lemma}[Subcritical Morrey potential bound]
\label{gt:lem-sobolev-morrey-potential-bound}\dcref{ch7:7.18}\group{ch07-sobolev-spaces}\source{gilbarg-trudinger:page-0177}
Let $f\in M^p(\Omega)$ and $\delta=p^{-1}<\mu$. Then
\begin{equation}
\tag{7.47}
|V_\mu f(x)|\leq
\frac{1-\delta}{\mu-\delta}
(\diam\Omega)^{n(\mu-\delta)}\|f\|_{M^p(\Omega)}
\quad\text{for almost every }x\in\Omega.
\end{equation}
\end{lemma}
\begin{proof}
\sketch Extend $f$ by zero and set
$v(\rho)=\int_{B_\rho(x)}|f|$. Stieltjes integration by parts in
$\int_0^d\rho^{n(\mu-1)}v'(\rho)\,d\rho$, where
$d=\diam\Omega$, followed by (7.46), gives (7.47).
\end{proof}

\begin{theorem}[Morrey growth criterion]
\label{gt:thm-sobolev-morrey-growth}\dcref{ch7:7.19}\group{ch07-sobolev-spaces}\source{gilbarg-trudinger:page-0177}
Let $u\in W^{1,1}(\Omega)$. Suppose $K>0$, $0<\alpha\leq1$, and
\begin{equation}
\tag{7.48}
\int_{B_R}|Du|\,dx\leq KR^{n-1+\alpha}
\end{equation}
for every ball $B_R\subset\Omega$. Then $u\in C^{0,\alpha}(\Omega)$ and
\begin{equation}
\tag{7.49}
\operatorname*{osc}_{B_R}u\leq C K R^\alpha,
\qquad C=C(n,\alpha).
\end{equation}
If $\Omega=\widetilde\Omega\cap\mathbb R_+^n$ and (7.48) holds for every
$B_R\subset\widetilde\Omega$, the same conclusion holds up to the flat
boundary, with (7.49) for those balls.
\end{theorem}
\begin{proof}
\sketch Apply Lemma~7.16 on balls or half-balls and then use Lemma~7.18 with
the Morrey exponent determined by (7.48).
\end{proof}

\begin{lemma}[Critical Morrey potential estimate]
\label{gt:lem-sobolev-morrey-potential-exponential}\dcref{ch7:7.20}\group{ch07-sobolev-spaces}\source{gilbarg-trudinger:page-0177, gilbarg-trudinger:page-0178}
Let $p>1$, $f\in M^p(\Omega)$, $\mu=p^{-1}$, and $g=V_\mu f$. There are
$c_1,c_2>0$, depending only on $n,p$, such that, for
$K=\|f\|_{M^p(\Omega)}$,
\begin{equation}
\tag{7.50}
\int_\Omega\exp\!\left(\frac{|g|}{c_1K}\right)dx
\leq c_2(\diam\Omega)^n.
\end{equation}
\end{lemma}
\begin{proof}
\sketch Factor the potential kernel and apply H\"older's inequality to bound $|g|$ by
powers of $V_{\mu/q}|f|$ and $V_{\mu+\mu/q}|f|$. Lemmas~7.12 and 7.18 give
$\int_\Omega|g|^q\leq C\{(p-1)qK\}^q(\diam\Omega)^n$. Summing moments proves
(7.50) when $c_1>(p-1)e$.
\end{proof}

\begin{theorem}[John--Nirenberg estimate]
\label{gt:thm-sobolev-john-nirenberg}\dcref{ch7:7.21}\group{ch07-sobolev-spaces}\source{gilbarg-trudinger:page-0178}
Let $\Omega$ be convex and $u\in W^{1,1}(\Omega)$. Suppose
\begin{equation}
\tag{7.51}
\int_{\Omega\cap B_R}|Du|\,dx\leq KR^{n-1}
\end{equation}
for every ball $B_R$. There are $\sigma_0,C>0$, depending only on $n$, such
that
\begin{equation}
\tag{7.52}
\int_\Omega\exp\!\left(\frac\sigma K|u-u_\Omega|\right)dx
\leq C(\diam\Omega)^n,
\qquad
\sigma=\sigma_0|\Omega|(\diam\Omega)^{-n}.
\end{equation}
\end{theorem}
\begin{proof}
\sketch Lemma~7.16 bounds $|u-u_\Omega|$ by a constant multiple of
$V_{1/n}|Du|$. Apply Lemma~7.20 with the Morrey bound (7.51).
\end{proof}

\section*{7.10. Compactness Results}

A continuous embedding $\mathcal B_1\hookrightarrow\mathcal B_2$ is compact
when bounded subsets of $\mathcal B_1$ have relatively compact images in
$\mathcal B_2$.

\begin{theorem}[Kondrachov compactness]
\label{gt:thm-sobolev-kondrachov}\dcref{ch7:7.22}\group{ch07-sobolev-spaces}\source{gilbarg-trudinger:page-0179, gilbarg-trudinger:page-0180}
If $p<n$, then $W_0^{1,p}(\Omega)$ embeds compactly into $L^q(\Omega)$ for
every $1\leq q<np/(n-p)$. If $p>n$, it embeds compactly into
$C^0(\overline\Omega)$.
\end{theorem}
\begin{proof}
\sketch For a bounded family in $W_0^{1,p}$, fixed-scale mollifications are
uniformly bounded and equicontinuous, hence precompact. The estimate
$\|u-u_h\|_1\leq h\|Du\|_1$ makes this family precompact in $L^1$.
Interpolation with Theorem~7.10 gives compactness in every subcritical
$L^q$. For $p>n$, use the uniform H\"older estimate of Theorem~7.17 and
Arzela--Ascoli.
\end{proof}

\begin{corollary}[Higher-order compact embeddings]
\label{gt:cor-sobolev-kondrachov-higher-order}
\dcref{ch7:7.22}
\group{ch07-sobolev-spaces}
\source{gilbarg-trudinger:page-0180}
The embeddings
\[
W_0^{k,p}(\Omega)\hookrightarrow L^q(\Omega)
\quad\left(kp<n,\ q<\frac{np}{n-kp}\right)
\]
and
\[
W_0^{k,p}(\Omega)\hookrightarrow C^m(\overline\Omega)
\quad\left(0\leq m<k-\frac np\right)
\]
are compact.
\end{corollary}
\begin{proof}
\sketch Iterate Theorem~7.22 through the derivatives of order below $k$.
\end{proof}

\section*{7.11. Difference Quotients}

\begin{definition}[Difference quotient]
\label{gt:def-sobolev-difference-quotient}
\dcref{ch7:7.53}
\group{ch07-sobolev-spaces}
\source{gilbarg-trudinger:page-0180}
For the $i$th coordinate vector $e_i$ and $h\neq0$, set
\begin{equation}
\tag{7.53}
\Delta_i^h u(x)=\Delta^h u(x)
=\frac{u(x+he_i)-u(x)}h.
\end{equation}
\end{definition}

\begin{lemma}[Difference quotients of Sobolev functions]
\label{gt:lem-sobolev-difference-quotient-bound}\dcref{ch7:7.23}\group{ch07-sobolev-spaces}\source{gilbarg-trudinger:page-0180, gilbarg-trudinger:page-0181}
If $u\in W^{1,p}(\Omega)$ and $\Omega'\Subset\Omega$ with
$|h|<\dist(\Omega',\partial\Omega)$, then
$\Delta_i^hu\in L^p(\Omega')$ and
\[
\|\Delta_i^hu\|_{L^p(\Omega')}\leq\|D_i u\|_{L^p(\Omega)}.
\]
\end{lemma}
\begin{proof}
\sketch For smooth $u$, write the difference quotient as the average of
$D_i u$ along a segment, apply H\"older's inequality, and integrate over $\Omega'$. Extend
the estimate by Theorem~7.9.
\end{proof}

\begin{lemma}[Difference quotient criterion]
\label{gt:lem-sobolev-difference-quotient-criterion}\dcref{ch7:7.24}\group{ch07-sobolev-spaces}\source{gilbarg-trudinger:page-0181}
Let $1<p<\infty$ and $u\in L^p(\Omega)$. Suppose a constant $K$ bounds
$\|\Delta_i^hu\|_{L^p(\Omega')}$ for every $\Omega'\Subset\Omega$ and every
$h$ with $0<|h|<\dist(\Omega',\partial\Omega)$. Then $D_i u$ exists weakly
and $\|D_i u\|_{L^p(\Omega)}\leq K$.
\end{lemma}
\begin{proof}
\sketch Weak compactness gives a sequence $h_m\to0$ for which
$\Delta_i^{h_m}u$ converges weakly to some $v$ with $\|v\|_p\leq K$. The
discrete integration-by-parts identity
$\int\varphi\Delta_i^hu=-\int u\Delta_i^{-h}\varphi$ identifies
$v=D_i u$.
\end{proof}

\section*{7.12. Extension and Interpolation}

\begin{theorem}[Sobolev extension]
\label{gt:thm-sobolev-extension}\dcref{ch7:7.25}\group{ch07-sobolev-spaces}\source{gilbarg-trudinger:page-0182, gilbarg-trudinger:page-0183}
Let $k\geq1$, $1\leq p<\infty$, and let $\Omega$ be a
$C^{k-1,1}$ domain. Then $C^\infty(\overline\Omega)$ is dense in
$W^{k,p}(\Omega)$. For every open $\Omega'$ with
$\overline\Omega\subset\Omega'$, there is a
bounded linear operator
$E:W^{k,p}(\Omega)\to W_0^{k,p}(\Omega')$ such that $Eu=u$ on $\Omega$ and
\begin{equation}
\tag{7.54}
\|Eu\|_{k,p;\Omega'}\leq C\|u\|_{k,p;\Omega},
\qquad C=C(k,\Omega,\Omega').
\end{equation}
\end{theorem}
\begin{proof}
\sketch On the half-space, the translated mollifications
\begin{equation}
\tag{7.55}
v_h(x)=u_h(x+2he_n)
\end{equation}
converge to $u$ in $W^{k,p}(\mathbb R_+^n)$. Define the reflection extension
\begin{equation}
\tag{7.56}
E_0u(x)=
\begin{cases}
u(x),&x_n>0,\\
\displaystyle\sum_{i=1}^k c_i u(x',-x_n/i),&x_n<0,
\end{cases}
\end{equation}
where $\sum_{i=1}^kc_i(-1/i)^m=1$ for $0\leq m<k$. It satisfies
\begin{equation}
\tag{7.57}
\|E_0u\|_{k,p;\mathbb R^n}
\leq C(k)\|u\|_{k,p;\mathbb R_+^n}.
\end{equation}
Flatten the boundary by finitely many $C^{k-1,1}$ charts $\psi_j$ and use a
subordinate partition of unity $\eta_j$. With an interior term $u\eta_0$,
the global extension is
\begin{equation}
\tag{7.58}
Eu=u\eta_0+\sum_{j=1}^N
E_0[(\eta_ju)\circ\psi_j^{-1}]\circ\psi_j.
\end{equation}
The chart bounds and (7.57) imply (7.54); mollifying $Eu$ proves density.
\end{proof}

\begin{theorem}[General Sobolev embeddings]
\label{gt:thm-sobolev-general-embedding}\dcref{ch7:7.26}\group{ch07-sobolev-spaces}\source{gilbarg-trudinger:page-0183}
Let $\Omega$ be a $C^{0,1}$ domain.
\begin{enumerate}
\item If $kp<n$ and $p^*=np/(n-kp)$, then
$W^{k,p}(\Omega)\hookrightarrow L^{p^*}(\Omega)$ continuously and
$W^{k,p}(\Omega)\hookrightarrow L^q(\Omega)$ compactly for every $q<p^*$.
\item If $0\leq m<k-n/p<m+1$ and $\alpha=k-n/p-m$, then
$W^{k,p}(\Omega)\hookrightarrow C^{m,\alpha}(\overline\Omega)$ continuously
and into $C^{m,\beta}(\overline\Omega)$ compactly for every
$0\leq\beta<\alpha$.
\end{enumerate}
\end{theorem}
\begin{proof}
\sketch Apply the case $k=1$ of Theorem~7.25 on the Lipschitz domain, then
iterate the first-order embeddings and compactness results through the weak
derivatives of $u$.
\end{proof}

\begin{theorem}[Interpolation on $W_0^{k,p}$]
\label{gt:thm-sobolev-interpolation-zero}\dcref{ch7:7.27}\group{ch07-sobolev-spaces}\source{gilbarg-trudinger:page-0183, gilbarg-trudinger:page-0184}
If $u\in W_0^{k,p}(\Omega)$, $\varepsilon>0$, and $0<|\beta|<k$, then
\begin{equation}
\tag{7.59}
\|D^\beta u\|_{p;\Omega}
\leq\varepsilon\|u\|_{k,p;\Omega}
+C\varepsilon^{|\beta|/(|\beta|-k)}\|u\|_{p;\Omega},
\qquad C=C(k).
\end{equation}
\end{theorem}
\begin{proof}
\sketch For $u\in C_0^2(\mathbb R)$, apply the mean value theorem on
intervals of length $\varepsilon$, integrate over the choices of endpoints,
and use H\"older's inequality to obtain
\begin{equation}
\tag{7.60}
\int|u'|^p\leq2^{p-1}\left\{
\varepsilon^p\int|u''|^p+
\left(\frac{18}{\varepsilon}\right)^p\int|u|^p\right\}.
\end{equation}
Apply this estimate along coordinate lines, use density, and induct on
$k$ and $|\beta|$.
\end{proof}

\begin{theorem}[Interpolation on a domain]
\label{gt:thm-sobolev-interpolation-domain}\dcref{ch7:7.28}\group{ch07-sobolev-spaces}\source{gilbarg-trudinger:page-0185}
Let $\Omega$ be a $C^{1,1}$ domain, $u\in W^{k,p}(\Omega)$,
$\varepsilon>0$, and $0<|\beta|<k$. Then
\begin{equation}
\tag{7.61}
\|D^\beta u\|_{p;\Omega}
\leq\varepsilon\|u\|_{k,p;\Omega}
+C\varepsilon^{|\beta|/(|\beta|-k)}\|u\|_{p;\Omega},
\qquad C=C(k,\Omega).
\end{equation}
\end{theorem}
\begin{proof}
\sketch Extend $u$ by Theorem~7.25, apply Theorem~7.27 to the extension, and
use the extension estimate (7.54).
\end{proof}
